\documentclass[10pt,a4paper]{amsart}
\usepackage{amsmath,amssymb,amsthm,graphicx}
% For commutative diagrams we use Xy-pic
\usepackage[all]{xy}

% We use 2cell for 2-commutative diagrams.
\xyoption{2cell}
% Theorem environments.
%
\theoremstyle{plain}
\newtheorem{theorem}[subsection]{Theorem}
\newtheorem{proposition}[subsection]{Proposition}
\newtheorem{lemma}[subsection]{Lemma}

\theoremstyle{definition}
\newtheorem{definition}[subsection]{Definition}
\newtheorem{example}[subsection]{Example}
\newtheorem{exercise}[subsection]{Exercise}
\newtheorem{situation}[subsection]{Situation}

\theoremstyle{remark}
\newtheorem{remark}[subsection]{Remark}
\newtheorem{remarks}[subsection]{Remarks}

\numberwithin{equation}{subsection}

% Macros
%
\def\lim{\mathop{\mathrm{lim}}\nolimits}
\def\colim{\mathop{\mathrm{colim}}\nolimits}
\def\Spec{\mathop{\mathrm{Spec}}}
\def\Hom{\mathop{\mathrm{Hom}}\nolimits}
\def\Ext{\mathop{\mathrm{Ext}}\nolimits}
\def\SheafHom{\mathop{\mathcal{H}\!\mathit{om}}\nolimits}
\def\SheafExt{\mathop{\mathcal{E}\!\mathit{xt}}\nolimits}
\def\Sch{\mathit{Sch}}
\def\Mor{\mathop{\mathrm{Mor}}\nolimits}
\def\Ob{\mathop{\mathrm{Ob}}\nolimits}
\def\Sh{\mathop{\mathit{Sh}}\nolimits}
\def\NL{\mathop{N\!L}\nolimits}
\def\CH{\mathop{\mathrm{CH}}\nolimits}
\def\proetale{{pro\text{-}\acute{e}tale}}
\def\etale{{\acute{e}tale}}
\def\QCoh{\mathit{QCoh}}
\def\Ker{\mathop{\mathrm{Ker}}}
\def\Im{\mathop{\mathrm{Im}}}
\def\Coker{\mathop{\mathrm{Coker}}}
\def\Coim{\mathop{\mathrm{Coim}}}

% Boxtimes
%
\DeclareMathSymbol{\boxtimes}{\mathbin}{AMSa}{"02}

%
% Macros for moduli stacks/spaces
%
\def\QCohstack{\mathcal{QC}\!\mathit{oh}}
\def\Cohstack{\mathcal{C}\!\mathit{oh}}
\def\Spacesstack{\mathcal{S}\!\mathit{paces}}
\def\Quotfunctor{\mathrm{Quot}}
\def\Hilbfunctor{\mathrm{Hilb}}
\def\Curvesstack{\mathcal{C}\!\mathit{urves}}
\def\Polarizedstack{\mathcal{P}\!\mathit{olarized}}
\def\Complexesstack{\mathcal{C}\!\mathit{omplexes}}
% \Pic is the operator that assigns to X its picard group, usage \Pic(X)
% \Picardstack_{X/B} denotes the Picard stack of X over B
% \Picardfunctor_{X/B} denotes the Picard functor of X over B
\def\Pic{\mathop{\mathrm{Pic}}\nolimits}
\def\Picardstack{\mathcal{P}\!\mathit{ic}}
\def\Picardfunctor{\mathrm{Pic}}
\def\Deformationcategory{\mathcal{D}\!\mathit{ef}}
\usepackage{cleveref}
\crefname{originallemma}{original lemma}{original lemmas}
\Crefname{originallemma}{Original lemma}{Original lemmas}
\makeatletter
\newlabel{lemma-c1-cap-additive}{{02SP}{source}{}{sourcepointer-1}{}}
\newlabel{lemma-c1-cap-additive@cref}{{[originallemma][0][]02SP}{[1][0][]source}{}{}{}}
\makeatother
\crefname{originallemma}{original lemma}{original lemmas}
\Crefname{originallemma}{Original lemma}{Original lemmas}
\makeatletter
\newlabel{lemma-decompose-section}{{0F97}{source}{}{sourcepointer-2}{}}
\newlabel{lemma-decompose-section@cref}{{[originallemma][0][]0F97}{[1][0][]source}{}{}{}}
\makeatother
\crefname{originallemma}{original lemma}{original lemmas}
\Crefname{originallemma}{Original lemma}{Original lemmas}
\makeatletter
\newlabel{lemma-flat-pullback-cap-c1}{{02SS}{source}{}{sourcepointer-3}{}}
\newlabel{lemma-flat-pullback-cap-c1@cref}{{[originallemma][0][]02SS}{[1][0][]source}{}{}{}}
\makeatother
\crefname{originallemma}{original lemma}{original lemmas}
\Crefname{originallemma}{Original lemma}{Original lemmas}
\makeatletter
\newlabel{lemma-linebundle-formulae}{{0F96}{source}{}{sourcepointer-4}{}}
\newlabel{lemma-linebundle-formulae@cref}{{[originallemma][0][]0F96}{[1][0][]source}{}{}{}}
\makeatother
\crefname{originallemma}{original lemma}{original lemmas}
\Crefname{originallemma}{Original lemma}{Original lemmas}
\makeatletter
\newlabel{situation-setup}{{02QL}{source}{}{sourcepointer-5}{}}
\newlabel{situation-setup@cref}{{[originallemma][0][]02QL}{[1][0][]source}{}{}{}}
\makeatother
\title{Compatibility of the Gysin map with the two sections of the decomposed zero scheme}
\author{The Stacks Project Authors}
\date{stacks-project:chow.tex:local-snapshot-20261012; GNU FDL 1.2 or later}
\begin{document}
\maketitle
\noindent\textit{Editorial note.} All delivered mathematical text is the author's own, copied from the pinned Stacks Project chapter file chow.tex; no statement or proof step was rewritten, reordered or removed.\par\medskip
\noindent\textit{Editorial note.} Preamble provenance: the pinned chapter file opens with \textbackslash{}input\{preamble\} and carries no preamble text of its own, so the project's single Stacks-wide preamble staged at sources/stacks-chow/preamble.tex is excerpted verbatim as two spans, lines 14-18 (Xy-pic setup) and lines 154-217 (the native theorem declarations and mathematical macros). Nothing was added, removed or redefined; the wrapper class and the amsmath/amssymb/amsthm/graphicx base packages are supplied by the assembler.\par\medskip
\noindent\textit{Editorial note.} Same-chapter pointers are located at their \textbackslash{}label lines inside the pinned chapter file; cross-chapter pointers are located in the project's own tag list staged in the snapshot (sources/stacks-chow/tags.txt), at the verbatim lines carrying their canonical tags.\par\medskip
\noindent\textit{Allowed standard prerequisites (editorial).}\begin{itemize}
\item The decomposition of the zero scheme of a section of a tensor product of invertible sheaves. Same chapter chow.tex, label lemma-decompose-section, invoked in the statement at proof line 5885 and used throughout. Supplies the commutative diagram with p : L -> Y, the two sections s\_i, the line bundles D\_i -> D and the identification of D\_1 with the relative spectrum of the inverse of L\_2; the compatibility computation itself is in the counted span.
\item Additivity of capping with a first Chern class and compatibility with flat pullback. Same chapter chow.tex, labels lemma-c1-cap-additive (proof line 5900) and lemma-flat-pullback-cap-c1 (proof line 5925). Gives i\textasciicircum{}*[W] = c\_1(L\_1) cap [W] + c\_1(L\_2) cap [W] when W is contained in D, and allows c\_1 to be moved across the flat pullback p\_1\textasciicircum{}* in the chain of equalities.
\item Formulae for a line bundle and its zero section. Same chapter chow.tex, label lemma-linebundle-formulae, invoked at proof line 5927. Identifies the class of the zero section W\_1 intersect W\_2 of the line bundle W\_1 -> W with p\_1\textasciicircum{}*(c\_1(L\_2) cap [W]), which is the step that closes the first case.
\end{itemize}\medskip
\ifdefined\hypertarget\hypertarget{sourcepointer-1}{}\fi
\noindent\textit{Source reference.} Original lemma 02SP: sources/stacks-chow/pinned-original-source.tex, source lines 4330--4330. Navigation pointer for \textbackslash{}ref\{lemma-c1-cap-additive\} used inside the counted span (line 5900); the argument that consumes it is retained in the span.\par\medskip
\ifdefined\hypertarget\hypertarget{sourcepointer-2}{}\fi
\noindent\textit{Source reference.} Original lemma 0F97: sources/stacks-chow/pinned-original-source.tex, source lines 5822--5822. Navigation pointer for \textbackslash{}ref\{lemma-decompose-section\} used inside the counted span (line 5885); the argument that consumes it is retained in the span.\par\medskip
\ifdefined\hypertarget\hypertarget{sourcepointer-3}{}\fi
\noindent\textit{Source reference.} Original lemma 02SS: sources/stacks-chow/pinned-original-source.tex, source lines 4533--4533. Navigation pointer for \textbackslash{}ref\{lemma-flat-pullback-cap-c1\} used inside the counted span (line 5925); the argument that consumes it is retained in the span.\par\medskip
\ifdefined\hypertarget\hypertarget{sourcepointer-4}{}\fi
\noindent\textit{Source reference.} Original lemma 0F96: sources/stacks-chow/pinned-original-source.tex, source lines 5775--5775. Navigation pointer for \textbackslash{}ref\{lemma-linebundle-formulae\} used inside the counted span (line 5927); the argument that consumes it is retained in the span.\par\medskip
\ifdefined\hypertarget\hypertarget{sourcepointer-5}{}\fi
\noindent\textit{Source reference.} Original lemma 02QL: sources/stacks-chow/pinned-original-source.tex, source lines 1372--1372. Navigation pointer for \textbackslash{}ref\{situation-setup\} used inside the counted span (line 5887); the argument that consumes it is retained in the span.\par\medskip
\begin{lemma}
\label{lemma-decompose-section-formulae}
In the situation of Lemma \ref{lemma-decompose-section}
assume $Y$ is locally of finite type over $(S, \delta)$ as in
Situation \ref{situation-setup}. Then we have
$i_1^*p^*\alpha = p_1^*i^*\alpha$
in $\CH_k(D_1)$ for all $\alpha \in \CH_k(Y)$.
\end{lemma}

\begin{proof}
Let $W \subset Y$ be an integral closed subscheme of $\delta$-dimension $k$.
We distinguish two cases.

\medskip\noindent
Assume $W \subset D$. Then
$i^*[W] = c_1(\mathcal{L}_1) \cap [W] + c_1(\mathcal{L}_2) \cap [W]$
in $\CH_{k - 1}(D)$ by our definition of Gysin homomorphisms and the
additivity of Lemma \ref{lemma-c1-cap-additive}.
Hence $p_1^*i^*[W] =
p_1^*(c_1(\mathcal{L}_1) \cap [W]) + p_1^*(c_1(\mathcal{L}_2) \cap [W])$.
On the other hand, we have
$p^*[W] = [p^{-1}(W)]_{k + 1}$ by construction of flat pullback.
And $p^{-1}(W) = W_1 \cup W_2$ (scheme theoretically)
where $W_i = p_i^{-1}(W)$ is a line bundle over $W$
by the lemma (since formation of the diagram commutes with base change).
Then $[p^{-1}(W)]_{k + 1} = [W_1] + [W_2]$ as $W_i$ are integral closed
subschemes of $L$ of $\delta$-dimension $k + 1$. Hence
\begin{align*}
i_1^*p^*[W]
& =
i_1^*[p^{-1}(W)]_{k + 1} \\
& =
i_1^*([W_1] + [W_2]) \\
& =
c_1(p_1^*\mathcal{L}_1) \cap [W_1] + [W_1 \cap W_2]_k \\
& =
c_1(p_1^*\mathcal{L}_1) \cap p_1^*[W] + [W_1 \cap W_2]_k \\
& =
p_1^*(c_1(\mathcal{L}_1) \cap [W]) + [W_1 \cap W_2]_k
\end{align*}
by construction of Gysin homomorphisms, the definition of flat pullback
(for the second equality), and compatibility of $c_1 \cap -$
with flat pullback (Lemma \ref{lemma-flat-pullback-cap-c1}).
Since $W_1 \cap W_2$ is the zero section of the line bundle
$W_1 \to W$ we see from Lemma \ref{lemma-linebundle-formulae}
that $[W_1 \cap W_2]_k = p_1^*(c_1(\mathcal{L}_2) \cap [W])$.
Note that here we use the fact that $D_1$ is the line bundle
which is the relative spectrum of the inverse of $\mathcal{L}_2$.
Thus we get the same thing as before.

\medskip\noindent
Assume $W \not \subset D$. In this case, both $i_1^*p^*[W]$
and $p_1^*i^*[W]$ are represented by the $k - 1$ cycle associated
to the scheme theoretic inverse image of $W$ in $D_1$.
\end{proof}

\end{document}
